\documentclass{article}

% Only basic packages for text and images
\usepackage[margin=1in]{geometry}
\usepackage{graphicx}
\usepackage{float}

% Title page information
\title{Distributed Systems Project 2: Design, Implementation, and Evaluation of a Distributed Computing System}

\author{Madison Gage \\ Benjamin Niccum \\ University of Texas at Arlington \\ CSE 5306 - Distributed Systems \\ Spring 2026}

\date{\today}

\begin{document}

% Title page
\maketitle
\thispagestyle{empty}
\newpage

% Table of Contents
\tableofcontents
\newpage

\section{Introduction}

This report presents the design, implementation, and evaluation of a distributed system that addresses fundamental challenges in distributed computing. The project demonstrates practical experience with system architecture trade-offs, communication models, performance optimization, and scalability considerations.

\subsection{Project Overview}
% Provide a brief overview of what the system does and its purpose

This project presents a distributed arithmetic processing system designed to demonstrate key principles of distributed computing through practical implementation. The system processes mathematical operations (addition, subtraction, multiplication, and division) across a network of containerized nodes, providing both interactive and batch processing capabilities.

\textbf{System Capabilities:} The arithmetic machine accepts computational requests through two primary input methods: manual entry via an interactive command-line interface, or automated batch processing from CSV files containing multiple operations. This dual-input approach enables evaluation of system performance under both lightweight interactive workloads and high-throughput batch scenarios.

\textbf{Deployment Architecture:} The system is fully containerized using Docker technology, requiring only a Docker installation on the target computing environment. Users can deploy the complete system by pulling the project from an online repository and executing the included build scripts, which utilize Docker Compose configurations to orchestrate the multi-node deployment.

\textbf{Distributed Processing Model:} The system employs a six-node distributed architecture comprising:
\begin{itemize}
    \item One user interface node responsible for input validation and request routing
    \item Four specialized computation nodes, each dedicated to a specific arithmetic operation type
    \item One display node that aggregates results, performs statistical analysis, and manages output
\end{itemize}

\textbf{Workflow Process:} When arithmetic operations are submitted, the user interface node validates input data and intelligently routes requests to the appropriate computation node based on the operation type. Following computation, results are forwarded to the display node, which stores the data persistently, generates computational statistics, and presents formatted output to the user. This architecture demonstrates load distribution, specialized processing, and centralized result aggregation patterns common in distributed systems.


\subsection{Objectives}
The primary objectives of this project include:
\begin{itemize}
    \item Design and implement a distributed system supporting six functional requirements
    \item Utilize microservice/object-based/resource-based architecture patterns
    \item Implement MPI as the primary communication model
    \item Deploy the system across six containerized nodes using Docker
    \item Evaluate performance and scalability trade-offs
    \item Document the role of AI tools in the development process
\end{itemize}

\subsection{Report Structure}
This report is organized as follows: the system architecture and design decisions are described first, followed by the six functional requirements, implementation details including containerization and MPI communication, evaluation methodology, experimental results, system design trade-offs, AI tool usage, and conclusions.

\section{System Design and Architecture}

\subsection{Architecture Overview}
% Describe your chosen architecture (microservice/object-based/resource-based)
% Include system architecture diagram

\begin{figure}[H]
    \centering
    \includegraphics[width=0.8\textwidth]{images/system_architecture.png}
    \caption{System Architecture Overview}
\end{figure}

We employed two prominent design principles during our project planning phase that guided the overall system architecture and implementation approach.

\textbf{Layered Communication Architecture:} The first design principle follows a simple layered approach for inter-node communication. The architecture consists of three distinct layers:

The top layer is the user interface, which communicates directly with users to gather computation requests. Once all operations are collected, the user interface sends messages to the second layer—the computation layer. The computation layer consists of four identical nodes running the same code, but each node is specialized to handle only one specific operation type (addition, subtraction, multiplication, or division).

The user interface layer primarily communicates with the computation layer throughout the processing phase. However, there is one important exception: a single termination message is sent directly to the third layer—the display layer—after all computations have been dispatched. This termination signal includes the expected number of computations and ensures proper synchronization. The display layer compares the expected count with the actual received count, and when they match, it sends termination signals back to the computation nodes, allowing them to safely end their processes. Finally, the display node completes its statistical analysis before terminating its own program.

\textbf{Object-Based Node Design:} The second design principle zooms in closer to the code level and follows an object-based approach. Each node functions as an object with one primary attribute—its running state—and several associated functions. Initially, all nodes start in a "not running" state, which toggles to "running" when the user activates the system. The nodes toggle back to "not running" when they receive a termination signal.

The core data object in our system is a struct called ComputationResult that gets passed between nodes throughout the entire process. For each computation that runs through the system, three instances of this struct are created, and the information is transferred from node to node using MPI functions. The ComputationResult struct contains seven key attributes: operation\_type, operand1, operand2, result, request\_id, status, and time. These attributes are updated and utilized by different functions as the computation progresses through the distributed system.


\subsection{Component Interaction}

The distributed arithmetic system operates through coordinated interactions between six specialized nodes, each fulfilling distinct roles in the computational workflow. This section outlines the high-level communication flows, data lifecycle, and coordination mechanisms that enable distributed processing.

\textbf{System Workflow Overview:}

The computational workflow follows a structured three-phase pattern: request initiation, distributed processing, and result aggregation. Users interact exclusively with the user interface node, which serves as both the entry point and coordinator for all computational activities. The four computation nodes operate as specialized processors, while the display node manages result collection and system output.

\textbf{Request Distribution Phase:}
The user interface node accepts computational requests through dual input mechanisms and performs critical preprocessing functions including input validation, data structure creation, and intelligent request routing. Based on operation type, requests are distributed to specialized computation nodes: addition (node 1), subtraction (node 2), multiplication (node 3), and division (node 4). This specialization enables parallel processing of different operation types while maintaining focused computational expertise.

\textbf{Distributed Processing Phase:}
Computation nodes operate independently and asynchronously, processing incoming requests without coordination overhead between nodes. Each node performs arithmetic calculations, timing measurements, and error detection before generating structured results. The absence of inter-computation-node communication maximizes processing efficiency and eliminates potential synchronization bottlenecks.

\textbf{Result Aggregation Phase:}
The display node serves as the centralized collector for all computation results, implementing sophisticated tracking mechanisms to ensure completeness and consistency. It maintains comprehensive statistics, performs result validation, generates output files, and coordinates orderly system termination. The aggregation process includes real-time statistical analysis and persistent storage of all computational results.

\textbf{Data Lifecycle Management:}
The ComputationResult structure serves as the primary data container throughout the system lifecycle. Initially created and populated in the user interface, the structure travels through computation nodes where results and timing data are added, finally arriving at the display node for aggregation and analysis. This consistent data format ensures reliable information flow across all system components.

\textbf{Coordination and Synchronization:}
System coordination relies on count-based validation and structured termination protocols. The user interface communicates expected computation counts to the display node, enabling validation of result completeness. Once all expected results are received, the display node initiates coordinated shutdown of all computation nodes, ensuring graceful system termination without data loss.

\textbf{Error Management:}
Error handling is integrated throughout the system workflow, with computation nodes detecting and reporting errors through standardized status codes. The display node processes both successful results and error conditions uniformly, maintaining system stability and providing comprehensive error reporting to users.



\subsection{Communication Model}
This section provides detailed technical specifications for the MPI-based communication infrastructure that enables reliable message passing across the distributed arithmetic system.

\textbf{MPI Implementation Architecture:}

The communication implementation utilizes core MPI functions optimized for blocking point-to-point communication. MPI\_Send() provides synchronous message transmission with guaranteed delivery confirmation, while MPI\_Recv() enables flexible message reception with source address wildcarding through MPI\_ANY\_SOURCE parameters. The MPI\_COMM\_WORLD communicator establishes a unified communication domain encompassing all six system nodes, identified through rank-based addressing (0-5).

Process initialization and cleanup utilize MPI\_Init() and MPI\_Finalize() to establish proper communication environment management, ensuring clean resource allocation and deallocation across distributed nodes. The blocking communication model eliminates buffer management complexity while providing strong consistency guarantees for message ordering and delivery.

\textbf{Protocol Specification and Message Tags:}

The message classification system implements four distinct tag protocols optimized for semantic message routing:

\begin{itemize}
\item \texttt{ALU\_TAG (0x01)}: Request transmission from UI node to computation nodes
\item \texttt{RESULT\_TAG (0x02)}: Result transmission from computation nodes to display node  
\item \texttt{COUNT\_TAG (0x03)}: Expected count transmission for synchronization validation
\item \texttt{TERMINATE\_TAG (0x04)}: Shutdown coordination broadcast to computation nodes
\end{itemize}

Tag-based routing eliminates message parsing overhead by enabling immediate semantic interpretation at reception, while message multiplexing capabilities allow multiple message types to coexist on the same communication channels without interference.

\textbf{Binary Serialization and Data Format:}

Struct transmission utilizes MPI\_BYTE datatype specification for direct binary serialization, eliminating text conversion overhead and preserving complete floating-point precision. The fixed-structure layout ensures consistent memory alignment across heterogeneous architectures:

\begin{verbatim}
struct ComputationResult {
    int operation_type;    // 4 bytes: 1-4 operation codes
    double operand1;       // 8 bytes: IEEE 754 double precision
    double operand2;       // 8 bytes: IEEE 754 double precision  
    double result;         // 8 bytes: computed result value
    int request_id;        // 4 bytes: sequential identifier
    int status;            // 4 bytes: error/success codes
    double time;           // 8 bytes: execution timing data
} // Total: 44 bytes fixed structure
\end{verbatim}

Binary transmission provides deterministic serialization performance and eliminates platform-specific text encoding issues, while the fixed-size structure enables predictable bandwidth utilization and memory management.

\textbf{Network Performance Optimization:}

The communication protocol prioritizes low-latency transmission through direct MPI\_Send operations, avoiding buffered or asynchronous alternatives that introduce scheduling overhead. Blocking semantics prevent buffer overflow conditions and provide natural backpressure mechanisms that regulate computation node processing rates.

Memory efficiency is maximized through in-place struct transmission without additional copying operations, while the binary format minimizes network bandwidth requirements compared to text-based protocols. Tag-based message dispatch enables zero-copy routing to appropriate handler functions upon reception.

\textbf{Reliability and Fault Detection:}

MPI's inherent reliability mechanisms provide guaranteed message delivery between functioning nodes, with automatic retransmission and error recovery handled transparently by the MPI implementation. Message ordering preservation ensures that sequential requests from the UI node arrive at computation nodes in transmission order, critical for maintaining request\_id correlation.

The communication layer implements robust error detection through MPI return code validation and embedded status field verification. Timeout protection is provided through MPI's blocking semantics, which prevent indefinite waiting conditions while maintaining system responsiveness.

\textbf{Scalability and Performance Characteristics:}

The point-to-point communication model scales linearly with node count, avoiding broadcast bottlenecks or centralized message routing overhead. Direct addressing through MPI ranks eliminates directory lookup requirements, while the tag-based protocol ensures O(1) message classification performance regardless of system size.

Bandwidth utilization remains constant per computation regardless of total system load, providing predictable performance characteristics for capacity planning. The communication protocol supports heterogeneous processing rates across computation nodes without requiring explicit load balancing mechanisms.


\section{Functional Requirements}

This section details the six functional requirements that our distributed system supports:

\subsection{Requirement 1: Distributed Input Processing}
\textbf{Description:} 
The system shall accept computational requests through both interactive CLI and batch CSV processing, with input validation and request queuing across distributed nodes.

\textbf{Implementation:} 
The user interface node implements dual input modes: an interactive command-line interface for manual operation entry and a CSV file parser for batch processing. Input validation occurs at the UI node level, checking for proper arithmetic format, valid operation types, and numeric operand formats. The system queues validated requests and distributes them to appropriate computation nodes based on operation type.

\subsection{Requirement 2: Parallel Arithmetic Operations}
\textbf{Description:} 
The system shall process addition, subtraction, multiplication, and division operations concurrently across specialized computation nodes with load balancing.

\textbf{Implementation:} 
Four dedicated computation nodes are deployed, each specialized for one arithmetic operation type. The UI node implements intelligent routing to distribute operations to the appropriate specialized node. Each computation node runs identical code but processes only its assigned operation type, enabling true parallel processing of different arithmetic operations simultaneously.

\subsection{Requirement 3: Fault-Tolerant Communication}
\textbf{Description:} 
The system shall implement reliable message passing between nodes using MPI with error detection and graceful handling of node failures.

\textbf{Implementation:} 
MPI message passing is implemented with structured communication protocols using defined message tags (COMPUTE_TAG, RESULT_TAG, TERMINATE_TAG). The system includes error status tracking within the ComputationResult struct, handles division-by-zero gracefully, and implements proper termination synchronization to ensure all nodes shut down cleanly.

\subsection{Requirement 4: Synchronous Result Aggregation}
\textbf{Description:} 
The system shall synchronize computation results from distributed nodes and maintain consistency through centralized result collection and validation.

\textbf{Implementation:} 
The display node serves as the centralized result aggregator, collecting all computation results from the four computation nodes. It implements result counting and validation by comparing expected computation counts with received results. The synchronization mechanism ensures all computations complete before statistical analysis begins and program termination occurs.

\subsection{Requirement 5: Performance Monitoring and Analytics}
\textbf{Description:} 
The system shall collect timing metrics, throughput statistics, and system performance data across all distributed components in real-time.

\textbf{Implementation:} 
Each computation node records execution timing using MPI_Wtime() functions, capturing precise computation duration. The display node aggregates timing data and generates statistical analytics including total computations processed, average computation times per operation type, and overall system throughput metrics. Results are stored persistently in CSV format for analysis.

\subsection{Requirement 6: Scalable Container Orchestration}
\textbf{Description:} 
The system shall support dynamic scaling and orchestration of containerized nodes with automated deployment and termination coordination.

\textbf{Implementation:} 
The system is fully containerized using Docker with Docker Compose orchestration managing the six-node deployment. Each node runs in its own container with defined inter-container communication via MPI. The deployment includes automated container startup sequencing, network configuration for MPI communication, and coordinated shutdown procedures ensuring all containers terminate gracefully.

\section{Implementation Details}

\subsection{Technology Stack}

The system utilizes a focused technology stack optimized for high-performance distributed computing:

\textbf{Programming Language:} C - Selected for its performance efficiency, low-level system control, and native MPI support. C provides optimal execution speed for arithmetic computations and minimal overhead for inter-process communication.

\textbf{Communication:} MPI (Message Passing Interface) - Chosen as the primary communication framework for its robust support of distributed computing paradigms, reliable message delivery, and proven scalability in high-performance computing environments.

\textbf{Containerization:} Docker - Employed to ensure consistent deployment environments across different systems, simplify dependency management, and enable portable distributed system deployment with minimal configuration requirements.

\textbf{Orchestration:} N/A - The system uses Docker Compose for container coordination rather than complex orchestration frameworks, keeping the deployment lightweight and manageable for the six-node architecture.

\textbf{Database:} N/A - Results are stored in CSV format for simplicity and portability, eliminating the need for database infrastructure while maintaining data persistence and analysis capabilities.

\textbf{Monitoring:} N/A - Performance monitoring is implemented through built-in MPI timing functions and custom analytics within the display node, providing sufficient metrics collection without external monitoring tools.

\subsection{MPI Implementation}
% Detail your MPI service definitions and implementation

\subsection{Dockerization Strategy}
% Explain how each node is containerized

\subsection{Node Architecture}
% Describe the architecture of individual nodes

\begin{figure}[H]
    \centering
    \includegraphics[width=0.6\textwidth]{images/node_architecture.png}
    \caption{Individual Node Architecture}
\end{figure}

\subsection{Deployment Configuration}
% Discuss Docker Compose and deployment strategies

\section{Evaluation Methodology}

\subsection{Evaluation Framework}
% Describe your approach to evaluating the system

\subsection{Comparison Systems}
This evaluation compares two system designs:

\subsubsection{Distributed Design}
% Describe your main distributed system design

\subsubsection{Centralized Design}
% Describe the alternative centralized design for comparison

\subsection{Performance Metrics}
The following metrics are used to evaluate system performance:

\begin{itemize}
    \item \textbf{Throughput:} Number of requests processed per second
    \item \textbf{Latency:} Response time for individual requests
    \item \textbf{Scalability:} Performance under varying workloads
    \item \textbf{Resource Utilization:} CPU, memory, and network usage
\end{itemize}

\subsection{Workload Specifications}
% Detail the different workloads used for testing

\section{Experimental Setup}

\subsection{Hardware Environment}
% Describe the hardware used for experiments

CPU: [CPU Details]

RAM: [Memory Details]

Storage: [Storage Details]

Network: [Network Details]

Operating System: [OS Details]

\subsection{Node Configuration}
% Explain how the six nodes are configured and deployed

\subsection{Test Environment Setup}
% Detail the setup process and configuration

\subsection{Measurement Tools}
% Describe tools used for performance measurement

\section{Results and Analysis}

\subsection{Performance Results}

\subsubsection{Throughput Analysis}

\begin{figure}[H]
    \centering
    \includegraphics[width=0.8\textwidth]{../images/throughput_comparison.png}
    \caption{Throughput Comparison: Distributed vs. Centralized Design}
\end{figure}

\subsubsection{Latency Analysis}

\begin{figure}[H]
    \centering
    \includegraphics[width=0.8\textwidth]{../images/latency_comparison.png}
    \caption{Latency Comparison: Distributed vs. Centralized Design}
\end{figure}

\subsection{Scalability Results}

\begin{figure}[H]
    \centering
    \includegraphics[width=0.8\textwidth]{../images/scalability_results.png}
    \caption{Scalability Analysis Under Varying Workloads}
\end{figure}

\subsection{Resource Utilization}

Distributed System:
- CPU Usage: [Data]\%
- Memory Usage: [Data]\%
- Network I/O: [Data] MB/s

Centralized System:
- CPU Usage: [Data]\%
- Memory Usage: [Data]\%
- Network I/O: [Data] MB/s

\subsection{Performance Analysis}
% Analyze and interpret the results

\section{System Design Trade-offs}

\subsection{Distributed vs. Centralized Architecture}

\subsubsection{Advantages of Distributed Design}
\begin{itemize}
    \item \textbf{Scalability:} 
    \item \textbf{Fault Tolerance:} 
    \item \textbf{Performance:} 
    \item \textbf{Resource Distribution:} 
\end{itemize}

\subsubsection{Disadvantages of Distributed Design}
\begin{itemize}
    \item \textbf{Complexity:} 
    \item \textbf{Network Overhead:} 
    \item \textbf{Consistency Challenges:} 
    \item \textbf{Debugging Difficulty:} 
\end{itemize}

\subsection{Communication Model Trade-offs}
% Discuss MPI advantages and limitations

\subsection{Containerization Benefits and Costs}
% Analyze Docker containerization impact

\subsection{Overall System Trade-offs}
% Summarize the key trade-offs identified

\section{AI Tools Usage and Lessons Learned}

\subsection{AI Tools Utilized}
% List and describe the AI tools used in the project

Tool 1: [Purpose] - Impact: [High/Medium/Low]

Tool 2: [Purpose] - Impact: [High/Medium/Low]

Tool 3: [Purpose] - Impact: [High/Medium/Low]

\subsection{Development Process Enhancement}

\subsubsection{Code Generation and Assistance}
% Describe how AI helped with code generation

\subsubsection{Debugging and Problem Solving}
% Explain AI assistance in debugging

\subsubsection{Documentation and Design}
% Discuss AI support in documentation

\subsection{Lessons Learned}

\subsubsection{Effective AI Tool Usage}
% What worked well with AI tools

\subsubsection{Limitations and Challenges}
% What limitations were encountered

\subsubsection{Best Practices}
% Recommendations for future AI tool usage

\subsection{Impact on Development Efficiency}
% Quantify the impact of AI tools on development time and quality

\section{Conclusion}

\subsection{Project Summary}
% Summarize the key achievements of the project

\subsection{Performance Insights}
% Summarize key performance findings

\subsection{Design Lessons}
% Key lessons learned about distributed system design

\subsection{Future Work}
% Potential improvements and extensions

\subsection{Final Reflections}
% Overall reflections on the project experience

% References
\newpage
\section{References}

[1] MPI Documentation. Available: https://MPI.io/

[2] Docker Documentation. Available: https://docs.docker.com/

[3] Add your references here as needed

\end{document}
